% !TeX root = book.tex
\section{Mục tiêu kiểm chứng thiết kế cơ sở dữ liệu}\label{s:valobjectives}

Mục tiêu của việc kiểm chứng thiết kế cơ sở dữ liệu là xác minh rằng mô hình quan hệ được đề xuất cung cấp cấu trúc cơ sở dữ liệu rõ ràng cho toàn bộ yêu cầu chức năng và ba mươi chín quy tắc nghiệp vụ đã xác định. Thay vì giả định thời gian thực thi hoặc kết quả đạt/không đạt trên dữ liệu tổng hợp khi chưa có môi trường sản xuất thực tế, quá trình kiểm chứng tập trung vào tính đúng đắn về cấu trúc, tính đầy đủ của quan hệ, khả năng thực thi ràng buộc, chuẩn hóa và khả năng truy vết của lược đồ.

Cụ thể, quá trình kiểm chứng đánh giá:
\begin{enumerate}
    \item \textbf{Mức độ bao phủ thực thể và bản số}: Bảo đảm mọi khái niệm vật lý và logic đều có quan hệ tương ứng, đồng thời bản số phản ánh chính xác thực tế vận hành.
    \item \textbf{Thực thi ràng buộc}: Xác định quy tắc nào được thực thi trực tiếp bằng cơ chế quan hệ (khóa chính, khóa ngoại, ràng buộc duy nhất, ràng buộc kiểm tra và chỉ mục từng phần), cũng như quy tắc nào cần trigger, giao dịch hoặc logic tầng ứng dụng.
    \item \textbf{Toàn vẹn tham chiếu và chính sách xóa}: Kiểm chứng hành vi xóa lan truyền, hạn chế xóa và đặt giá trị rỗng trong các quan hệ phụ thuộc cha--con.
    \item \textbf{Bảo toàn lịch sử và tách biệt kiểm toán}: Xác nhận các lần chuyển trạng thái nghiệp vụ và nhật ký kiểm toán bảo mật được tách thành những bảng chuyên biệt, ổn định.
    \item \textbf{Tính chặt chẽ của chuẩn hóa}: Chứng minh cả hai mươi ba quan hệ đều thỏa mãn Dạng chuẩn thứ ba (3NF), qua đó loại bỏ bất thường khi cập nhật, thêm và xóa dữ liệu.
    \item \textbf{Sự phù hợp giữa chỉ mục và truy vấn}: Xác minh chiến lược chỉ mục tương ứng trực tiếp với các mẫu truy cập truy vấn vận hành dự kiến.
\end{enumerate}

\section{Kiểm chứng các ràng buộc}\label{s:constraintval}

Các ràng buộc toàn vẹn là tuyến bảo vệ đầu tiên cho tính nhất quán dữ liệu. Lược đồ chuyển trực tiếp các yêu cầu nghiệp vụ thành những ràng buộc PostgreSQL tiêu chuẩn.

\subsection{Ràng buộc khóa chính}
Mỗi bảng trong hệ thống được gán một khóa chính thay thế gồm một cột do hệ quản trị cơ sở dữ liệu tự sinh (\texttt{BIGINT GENERATED ALWAYS AS IDENTITY}), hoặc một khóa chính ghép đối với các bảng liên kết.
\begin{itemize}
    \item \textbf{Bảng dữ liệu chủ và bảng vận hành}: \texttt{system\_user(user\_id)}, \texttt{customer(customer\_id)}, \texttt{customer\_address(address\_id)}, \texttt{delivery\_order(order\_id)}, \texttt{package(package\_id)}, \texttt{landing\_station(station\_id)}, \texttt{drone(drone\_id)}, \texttt{delivery\_activity(activity\_id)}, \texttt{tracking\_record(tracking\_id)}, \texttt{delivery\_confirmation(confirmation\_id)}.
    \item \textbf{Bảng liên kết}: \texttt{user\_role(user\_id, role\_id)} và \texttt{role\_permission(role\_id, permission\_id)} sử dụng khóa chính ghép, ngăn việc gán trùng vai trò hoặc quyền hạn.
\end{itemize}
\textbf{Kết luận kiểm chứng}: Được thực thi đầy đủ ở cấp cơ sở dữ liệu bằng chuỗi định danh và chỉ mục b-tree duy nhất của PostgreSQL.

\subsection{Ràng buộc khóa ngoại}
Các mối quan hệ giữa thực thể được thực thi thông qua khóa ngoại quan hệ.
\begin{itemize}
    \item Đơn hàng bắt buộc tham chiếu tới khách hàng hợp lệ (\texttt{delivery\_order.customer\_id $\to$ customer.customer\_id}).
    \item Hoạt động liên kết đơn hàng với máy bay (\texttt{delivery\_activity.order\_id $\to$ delivery\_order.order\_id}, \texttt{delivery\_activity.drone\_id $\to$ drone.drone\_id}).
    \item Dữ liệu đo từ xa và ngoại lệ liên kết với hoạt động giao hàng (\texttt{tracking\_record.activity\_id $\to$ delivery\_activity.activity\_id}, \texttt{delivery\_exception.activity\_id $\to$ delivery\_activity.activity\_id}).
    \item Hoạt động được lập lịch lại liên kết với lần thực hiện trước bằng khóa ngoại tự tham chiếu (\texttt{delivery\_activity.previous\_activity\_id $\to$ delivery\_activity.activity\_id}).
\end{itemize}
\textbf{Kết luận kiểm chứng}: Được thực thi đầy đủ ở cấp cơ sở dữ liệu. Các dòng tham chiếu tới bản ghi cha không tồn tại sẽ bị từ chối.

\subsection{Ràng buộc duy nhất}
Các khóa ứng viên và bản số một--một hoặc một--không/một được thực thi bằng ràng buộc \texttt{UNIQUE} ở cấp cơ sở dữ liệu:
\begin{itemize}
    \item \texttt{system\_user.username}, \texttt{system\_user.email}, \texttt{system\_user.phone\_number}.
    \item \texttt{customer.phone\_number}, \texttt{customer.email}.
    \item \texttt{drone.serial\_number}.
    \item \texttt{landing\_station.station\_code}.
    \item \texttt{delivery\_confirmation.order\_id} — Thực thi quy tắc một đơn giao hàng có tối đa một xác nhận thành công (BR-33).
    \item \texttt{delivery\_confirmation.activity\_id} — Bảo đảm mỗi lần thực hiện giao hàng có tối đa một bản ghi xác nhận (BR-35).
    \item \texttt{delivery\_route.activity\_id} — Bảo đảm một hoạt động có tối đa một tuyến đường dự kiến.
    \item \texttt{route\_point(route\_id, sequence\_number)} — Bảo đảm thứ tự điểm hành trình là duy nhất trong một tuyến đường.
\end{itemize}
\textbf{Kết luận kiểm chứng}: Được thực thi đầy đủ ở cấp cơ sở dữ liệu.

\subsection{Ràng buộc kiểm tra}
Tính toàn vẹn miền được kiểm chứng bằng các ràng buộc khai báo \texttt{CHECK}:
\begin{itemize}
    \item \textbf{Kích thước vật lý và tải trọng}: \texttt{package.weight\_kg > 0}, \texttt{package.length\_cm > 0}, \texttt{package.width\_cm > 0}, \texttt{package.height\_cm > 0}, \texttt{drone.max\_payload\_kg > 0}.
    \item \textbf{Tọa độ địa lý}: \texttt{latitude BETWEEN -90 AND 90}, \texttt{longitude BETWEEN -180 AND 180} trong \texttt{customer\_address}, \texttt{landing\_station}, \texttt{tracking\_record} và \texttt{route\_point}.
    \item \textbf{Miền giá trị vận hành}: \texttt{drone.battery\_level BETWEEN 0 AND 100}, \texttt{landing\_station.capacity >= 0}, \texttt{tracking\_record.altitude\_m >= 0}, \texttt{tracking\_record.speed\_kmh >= 0}.
    \item \textbf{Tập trạng thái rời rạc}: Được kiểm tra trên \texttt{delivery\_order.order\_status}, \texttt{package.package\_status}, \texttt{drone.status}, \texttt{landing\_station.status}, \texttt{delivery\_activity.activity\_status} và \texttt{package\_station\_event.event\_type}.
    \item \textbf{Giới hạn thời gian}: \texttt{delivery\_activity.started\_at <= delivery\_activity.ended\_at} và \texttt{station\_operator\_assignment.assigned\_from < station\_operator\_assignment.assigned\_to}.
\end{itemize}
\textbf{Kết luận kiểm chứng}: Được thực thi đầy đủ ở cấp cơ sở dữ liệu.

\section{Kiểm chứng từ quy tắc nghiệp vụ đến lược đồ}\label{s:brvalidation}

\textsc{Bảng}~\ref{tab:br_validation_part1} và \textsc{Bảng}~\ref{tab:br_validation_part2} trình bày kết quả kiểm chứng có hệ thống đối với toàn bộ ba mươi chín quy tắc nghiệp vụ, bao gồm cách biểu diễn quan hệ chính xác và mức thực thi cơ sở dữ liệu cụ thể.

\begin{table}[htbp]
\caption{Ma trận kiểm chứng: quy tắc về khách hàng, đơn hàng, máy bay và thực hiện giao hàng (BR-01 đến BR-24)}
\label{tab:br_validation_part1}
\centering
\small
\begin{tabular*}{\textwidth}{@{\extracolsep{\fill}} p{1.2cm} p{5.5cm} p{4.2cm} p{3.5cm}}
\toprule
\textbf{Mã} & \textbf{Tóm tắt quy tắc nghiệp vụ} & \textbf{Biểu diễn trong CSDL} & \textbf{Mức thực thi} \\
\midrule
BR-01 & Khách hàng có nhiều địa chỉ giao hàng & \texttt{customer 1:N customer\_address} & Được xác định bởi lược đồ \\
BR-02 & Đơn thuộc về một khách hàng & \texttt{delivery\_order.customer\_id} & CSDL thực thi (NOT NULL, FK) \\
BR-03 & Đơn sử dụng một địa chỉ giao hàng & \texttt{delivery\_order.delivery\_address\_id} & CSDL thực thi (NOT NULL, FK) \\
BR-04 & Địa chỉ thuộc về khách hàng đặt đơn & FK ghép \texttt{(delivery\_address\_id, customer\_id)} & CSDL thực thi (FK ghép) \\
BR-05 & Tối đa một địa chỉ mặc định & Chỉ mục từng phần \texttt{WHERE is\_default = TRUE} & CSDL thực thi (UNIQUE từng phần) \\
BR-06 & Đơn có ít nhất một kiện hàng & \texttt{package.order\_id FK} & Cần giao dịch/trigger \\
BR-07 & Đơn có thể chứa nhiều kiện hàng & \texttt{delivery\_order 1:N package} & Được xác định bởi lược đồ \\
BR-08 & Khối lượng, kích thước kiện hàng dương & \texttt{CHECK(weight\_kg > 0, etc.)} & CSDL thực thi (CHECK) \\
BR-09 & Vòng đời trạng thái kiện hàng & \texttt{CHECK(package\_status IN (...))} & CSDL thực thi (CHECK) \\
BR-10 & Một trạng thái đơn hiện tại & \texttt{delivery\_order.order\_status} & Được xác định bởi lược đồ \\
BR-11 & Bảo toàn thay đổi trạng thái đơn & \texttt{order\_status\_history} & Lược đồ (chính sách chỉ ghi nối tiếp) \\
BR-12 & Chuyển trạng thái hợp lệ & Mô hình chuyển trạng thái đơn & Cần trigger/ứng dụng \\
BR-13 & Đơn đã hủy không trở lại hoạt động & Mô hình chuyển trạng thái đơn & Cần trigger/ứng dụng \\
BR-14 & Đơn đã giao không được lập lịch lại & Mô hình chuyển trạng thái đơn & Cần trigger/ứng dụng \\
BR-15 & Máy bay phục vụ nhiều hoạt động & \texttt{drone 1:N delivery\_activity} & Được xác định bởi lược đồ \\
BR-16 & Không có chuyến bay đang hoạt động chồng lấn & Khoảng thời gian lịch bay & Cần trigger/ràng buộc loại trừ \\
BR-17 & Máy bay được phân công phải sẵn sàng & Kiểm tra trạng thái máy bay & Cần giao dịch/trigger \\
BR-18 & Khối lượng kiện $\le$ tải trọng máy bay & \texttt{SUM(weight) $\le$ max\_payload} & Cần trigger liên bản ghi/ứng dụng \\
BR-19 & Đơn có nhiều hoạt động giao hàng & \texttt{delivery\_order 1:N delivery\_activity} & Được xác định bởi lược đồ \\
BR-20 & Hoạt động là một lần thực hiện & Quan hệ \texttt{delivery\_activity} & Được xác định bởi lược đồ \\
BR-21 & Bảo toàn lần giao hàng thất bại & Dòng hoạt động được giữ (\texttt{'FAILED'}) & Chính sách/phân quyền \\
BR-22 & Lập lịch lại tạo hoạt động liên kết & \texttt{previous\_activity\_id} & CSDL thực thi (FK) \\
BR-23 & Hoạt động sử dụng chính xác một máy bay & \texttt{delivery\_activity.drone\_id} & CSDL thực thi (NOT NULL, FK) \\
BR-24 & Hoạt động có trạng thái/thời gian độc lập & \texttt{activity\_status}, các mốc thời gian & CSDL thực thi (NOT NULL, CHECK) \\
\bottomrule
\end{tabular*}
\end{table}

\begin{table}[htbp]
\caption{Ma trận kiểm chứng: quy tắc về trạm, xử lý kiện hàng, xác nhận và kiểm toán (BR-25 đến BR-39)}
\label{tab:br_validation_part2}
\centering
\small
\begin{tabular*}{\textwidth}{@{\extracolsep{\fill}} p{1.2cm} p{5.5cm} p{4.2cm} p{3.5cm}}
\toprule
\textbf{Mã} & \textbf{Tóm tắt quy tắc nghiệp vụ} & \textbf{Biểu diễn trong CSDL} & \textbf{Mức thực thi} \\
\midrule
BR-25 & Sức chứa trạm không âm & \texttt{CHECK(capacity >= 0)} & CSDL thực thi (CHECK) \\
BR-26 & Bảo toàn lịch sử trạng thái trạm & \texttt{station\_status\_history} & Lược đồ (chính sách chỉ ghi nối tiếp) \\
BR-27 & Khoảng thời gian bố trí nhân viên & \texttt{station\_operator\_assignment} & CSDL thực thi (CHECK, FK) \\
BR-28 & Kiểm tra quyền của nhân viên vận hành & Đối chiếu khoảng thời gian phân công & Cần trigger/ứng dụng \\
BR-29 & Kiện hàng được xử lý tại trạm & \texttt{package\_station\_event} & Được xác định bởi lược đồ \\
BR-30 & Ghi sự kiện lấy/bàn giao kiện hàng & \texttt{package\_station\_event.event\_type} & CSDL thực thi (NOT NULL, CHECK) \\
BR-31 & Sự kiện lưu trạm, nhân viên và thời gian & Các FK và mốc \texttt{occurred\_at} & CSDL thực thi (NOT NULL, FK) \\
BR-32 & Không bay khi chưa có sự kiện lấy hàng & Kiểm tra trình tự sự kiện & Cần trigger/ứng dụng \\
BR-33 & Tối đa một xác nhận cho mỗi đơn & \texttt{delivery\_confirmation.order\_id} & CSDL thực thi (UNIQUE) \\
BR-34 & Chỉ xác nhận khi đã giao hàng & Đối chiếu trạng thái đơn/hoạt động & Cần trigger/ứng dụng \\
BR-35 & Xác nhận liên kết hoạt động thành công & \texttt{delivery\_confirmation.activity\_id} & CSDL thực thi (UNIQUE, FK) \\
BR-36 & Bảo toàn người nhận, phương thức, bằng chứng & Các thuộc tính xác nhận & CSDL thực thi (NOT NULL) \\
BR-37 & Truy vết được các thao tác quan trọng & \texttt{audit\_log.user\_id}, \texttt{action\_time} & Được xác định bởi lược đồ \\
BR-38 & Kiểm toán lưu thay đổi trước/sau & \texttt{old\_value}, \texttt{new\_value} (\texttt{JSONB}) & Được xác định bởi lược đồ \\
BR-39 & Tách lịch sử nghiệp vụ và kiểm toán & Các bảng chuyên biệt độc lập & Tách biệt về cấu trúc \\
\bottomrule
\end{tabular*}
\end{table}

\section{Các cơ chế thực thi quan hệ nâng cao}\label{s:advancedenforcement}

Một số quy tắc nghiệp vụ là những thách thức điển hình trong thiết kế cơ sở dữ liệu quan hệ, vượt quá khả năng của các ràng buộc đơn cột thông thường. Thiết kế được đề xuất sử dụng các kỹ thuật mô hình hóa quan hệ nâng cao để trực tiếp thực thi những quy tắc này.

\subsection{Kiểm tra quyền sở hữu địa chỉ khách hàng (BR-04)}
Một điểm yếu thường gặp trong lược đồ đơn giản là cho phép đơn hàng chọn địa chỉ thuộc khách hàng khác nếu \texttt{customer\_id} và \texttt{delivery\_address\_id} chỉ được kiểm tra độc lập. Để ngăn vấn đề này ở cấp cơ sở dữ liệu:
\begin{enumerate}
    \item Định nghĩa ràng buộc duy nhất ghép trên \texttt{customer\_address(address\_id, customer\_id)}.
    \item Trong \texttt{delivery\_order}, thiết lập khóa ngoại ghép:
\begin{verbatim}
CONSTRAINT fk_order_address
  FOREIGN KEY (delivery_address_id, customer_id)
  REFERENCES customer_address(address_id, customer_id)
  ON DELETE RESTRICT
\end{verbatim}
\end{enumerate}
Cơ chế này bảo đảm hoàn toàn rằng không thể thêm một đơn hàng với địa chỉ thuộc khách hàng khác. Việc kiểm tra được thực thi trực tiếp bởi hệ quản trị cơ sở dữ liệu mà không cần mã ứng dụng hoặc trigger.

\subsection{Một địa chỉ mặc định cho mỗi khách hàng (BR-05)}
Nếu đặt ràng buộc \texttt{UNIQUE} tiêu chuẩn trên \texttt{is\_default}, toàn bộ bảng sẽ bị giới hạn sai thành chỉ có một địa chỉ mặc định. Vì vậy, thiết kế sử dụng tính năng chỉ mục từng phần của PostgreSQL:
\begin{verbatim}
CREATE UNIQUE INDEX uq_customer_default_address
  ON customer_address(customer_id)
  WHERE is_default = TRUE;
\end{verbatim}
Chỉ mục này chỉ đánh chỉ mục các dòng có \texttt{is\_default} bằng true. Mỗi khách hàng vẫn có thể có số lượng địa chỉ không mặc định tùy ý, trong khi quy tắc tối đa một địa chỉ mặc định được thực thi nghiêm ngặt ở cấp cơ sở dữ liệu.

\subsection{Bản số tối thiểu của kiện hàng (BR-06)}
Khóa ngoại tiêu chuẩn bảo đảm một kiện hàng tham chiếu một đơn ($N:1$), nhưng DDL quan hệ không thể tự thực thi việc mỗi đơn vừa được thêm phải có ngay ít nhất một kiện hàng ($1:1..N$ ở phía cha) trong một câu lệnh nguyên tử mà không tạo phụ thuộc vòng. Trong thiết kế này, đây là ràng buộc lược đồ cần được thực thi ở cấp giao dịch (thêm đơn và kiện hàng trong cùng một giao dịch cơ sở dữ liệu) hoặc bằng trigger ràng buộc trì hoãn.

\section{Toàn vẹn tham chiếu và chính sách xóa}\label{s:referentialval}

Lược đồ thiết lập các chính sách xóa (\texttt{ON DELETE}) rõ ràng và nhất quán cho toàn bộ khóa ngoại:

\begin{itemize}
    \item \textbf{Dữ liệu chủ phục vụ vận hành (\texttt{RESTRICT})}: Việc xóa dòng khỏi \texttt{customer}, \texttt{drone}, \texttt{landing\_station}, \texttt{delivery\_order} hoặc \texttt{delivery\_activity} bị chặn nếu tồn tại dòng con lịch sử hay vận hành phụ thuộc. Cơ chế này bảo vệ hồ sơ vận hành khỏi thao tác xóa lan truyền ngoài ý muốn.
    \item \textbf{Bảng liên kết bảo mật (\texttt{CASCADE})}: Khi xóa vai trò hoặc quyền hạn, thao tác xóa được lan truyền tới \texttt{user\_role} và \texttt{role\_permission}, vì các bản ghi liên kết không còn ý nghĩa khi thực thể bảo mật cha bị xóa.
    \item \textbf{Trách nhiệm người dùng và tham chiếu kiểm toán (\texttt{SET NULL})}: Khi tài khoản người dùng được lưu trữ hoặc loại bỏ, lịch sử vận hành và bản ghi kiểm toán vẫn được bảo toàn:
    \begin{itemize}
        \item \texttt{order\_status\_history.changed\_by} $\to$ \texttt{SET NULL}.
        \item \texttt{delivery\_activity.assigned\_by} $\to$ \texttt{SET NULL}.
        \item \texttt{delivery\_exception.resolved\_by} $\to$ \texttt{SET NULL}.
        \item \texttt{audit\_log.user\_id} $\to$ \texttt{SET NULL}.
    \end{itemize}
    Điều này bảo đảm dữ liệu đo từ xa, ngoại lệ và nhật ký kiểm toán lịch sử vẫn nguyên vẹn để đáp ứng yêu cầu tuân thủ, ngay cả khi bản ghi người dùng ban đầu bị thay đổi.
\end{itemize}

\section{Kiến trúc bảo toàn lịch sử và kiểm toán}\label{s:historyauditval}

Một yêu cầu thiết kế trung tâm là phân biệt lịch sử trạng thái nghiệp vụ với dấu vết kiểm toán bảo mật:

\begin{itemize}
    \item \textbf{Lịch sử nghiệp vụ (\texttt{order\_status\_history}, \texttt{station\_status\_history})}:
    \begin{itemize}
        \item Trả lời câu hỏi: \textit{``Trạng thái vận hành của thực thể nghiệp vụ này đã thay đổi như thế nào theo thời gian?''}
        \item Lưu các lần chuyển trạng thái nghiệp vụ rời rạc (\texttt{PENDING $\to$ APPROVED $\to$ ASSIGNED $\to$ IN\_TRANSIT $\to$ DELIVERED}) cùng thời điểm và lý do nghiệp vụ.
        \item Được thiết kế theo chính sách chỉ ghi nối tiếp ở cấp vận hành.
    \end{itemize}
    \item \textbf{Nhật ký kiểm toán (\texttt{audit\_log})}:
    \begin{itemize}
        \item Trả lời câu hỏi: \textit{``Ai đã sửa bản ghi nào, vào thời điểm nào và nội dung nào đã thay đổi?''}
        \item Lưu các thay đổi DML cấp thấp (\texttt{INSERT}, \texttt{UPDATE}, \texttt{DELETE}), ảnh chụp \texttt{JSONB} của bản ghi trước và sau thay đổi, địa chỉ IP máy khách và mã tác nhân.
        \item Sự tách biệt về cấu trúc bảo đảm nhật ký bảo mật có dung lượng lớn không gây nhiễu cho các truy vấn báo cáo nghiệp vụ.
    \end{itemize}
\end{itemize}

\section{Kiểm chứng chuẩn hóa}\label{s:normalizationval}

Thiết kế lược đồ được kiểm chứng chính thức theo các quy tắc chuẩn hóa quan hệ:

\subsection{Dạng chuẩn thứ nhất (1NF)}
\begin{itemize}
    \item \textbf{Tính nguyên tử}: Mọi miền giá trị của cột đều chứa các giá trị nguyên tử, không thể chia nhỏ. Đường phố, quận/huyện, thành phố và tọa độ được tách thành các cột riêng trong \texttt{customer\_address}.
    \item \textbf{Không có nhóm lặp}: Kích thước kiện hàng được lưu bằng các thuộc tính số vô hướng riêng (\texttt{length\_cm}, \texttt{width\_cm}, \texttt{height\_cm}) thay vì chuỗi văn bản phức hợp, chẳng hạn ``$30\times 20\times 10$''. Nhiều địa chỉ hoặc kiện hàng được mô hình hóa thành các dòng độc lập.
\end{itemize}

\subsection{Dạng chuẩn thứ hai (2NF)}
\begin{itemize}
    \item Mọi quan hệ đều ở 1NF.
    \item Tất cả thuộc tính không khóa phụ thuộc hàm đầy đủ vào khóa chính.
    \item Trong các bảng kết hợp có khóa ghép (\texttt{user\_role} với PK \texttt{(user\_id, role\_id)} và \texttt{role\_permission} với PK \texttt{(role\_id, permission\_id)}), không có thuộc tính không khóa nào phụ thuộc vào một tập con thực sự của khóa ghép. Trong \texttt{user\_role}, \texttt{assigned\_at} phụ thuộc vào cả người dùng và vai trò.
\end{itemize}

\subsection{Dạng chuẩn thứ ba (3NF)}
\begin{itemize}
    \item Mọi quan hệ đều ở 2NF.
    \item Không tồn tại phụ thuộc hàm bắc cầu ($X \to Y$ và $Y \to Z$, trong đó $Z$ là thuộc tính không khóa).
    \item \texttt{delivery\_order} chứa \texttt{customer\_id} dưới dạng khóa ngoại; tên, số điện thoại và thư điện tử của khách hàng không bị lặp trong \texttt{delivery\_order}.
    \item \texttt{delivery\_confirmation} tham chiếu \texttt{order\_id} và \texttt{activity\_id}; bảng không lưu mẫu máy bay hoặc thông tin liên hệ khách hàng.
    \item \texttt{tracking\_record} tham chiếu \texttt{activity\_id}; thông số máy bay được truy cập thông qua chuỗi quan hệ \texttt{tracking\_record $\to$ delivery\_activity $\to$ drone} mà không lặp dữ liệu.
\end{itemize}

\section{Kiểm chứng thiết kế chỉ mục và truy vấn}\label{s:indexval}

Các chỉ mục được thiết kế chặt chẽ để hỗ trợ những mẫu truy vấn vận hành dự kiến mà không tạo chi phí ghi không cần thiết trên các bảng có tần suất cao:

\begin{itemize}
    \item \textbf{Chỉ mục tra cứu khóa ngoại}: Cần thiết cho các phép nối quan hệ tần suất cao. Chỉ mục trên \texttt{delivery\_order(customer\_id)}, \texttt{order\_status\_history(order\_id)}, \texttt{delivery\_activity(order\_id)}, \texttt{delivery\_activity(drone\_id)}, \texttt{station\_operator\_assignment(station\_id)} và \texttt{package\_station\_event(package\_id)} ngăn quét tuần tự toàn bảng khi thực hiện phép nối.
    \item \textbf{Chỉ mục ghép cho chuỗi thời gian}: Bảng \texttt{tracking\_record} liên tục nhận các điểm GPS và dữ liệu đo từ xa trong chuyến bay. Chỉ mục ghép:
\begin{verbatim}
CREATE INDEX idx_tracking_activity_time
  ON tracking_record(activity_id, recorded_at DESC);
\end{verbatim}
    tối ưu việc truy xuất tọa độ đo từ xa gần nhất cho bảng điều khiển giám sát trực tiếp, đồng thời giới hạn phạm vi quét trong một hoạt động giao hàng.
    \item \textbf{Chỉ mục ràng buộc ngầm định}: Không tạo lại chỉ mục tường minh trên các cột đã được chỉ mục b-tree duy nhất do ràng buộc \texttt{UNIQUE} tạo ra (\texttt{system\_user.username}, \texttt{drone.serial\_number}, \texttt{landing\_station.station\_code}, \texttt{delivery\_confirmation.order\_id}).
\end{itemize}

\section{Giới hạn và ranh giới thực thi}\label{s:limitationsval}

Một thiết kế cơ sở dữ liệu chặt chẽ phải chỉ rõ ranh giới giữa nội dung có thể thực thi trực tiếp bằng định nghĩa lược đồ SQL khai báo và nội dung cần cơ chế thủ tục hoặc tầng ứng dụng:

\begin{enumerate}
    \item \textbf{Thực thi bằng lược đồ khai báo (trực tiếp trong PostgreSQL)}:
    \begin{itemize}
        \item Toàn vẹn thực thể (PK, cột định danh).
        \item Toàn vẹn tham chiếu (FK, chính sách xóa).
        \item Tính duy nhất (một địa chỉ mặc định bằng chỉ mục từng phần, một xác nhận cho mỗi đơn bằng ràng buộc duy nhất).
        \item Quyền sở hữu địa chỉ (FK ghép liên kết đơn và địa chỉ thông qua mã khách hàng).
        \item Kiểm tra miền vô hướng (khối lượng dương, tọa độ, mã trạng thái hợp lệ).
    \end{itemize}
    \item \textbf{Thực thi bằng thủ tục/trigger (cần logic được lưu trữ)}:
    \begin{itemize}
        \item \textbf{Giới hạn tải trọng liên bản ghi (BR-18)}: Việc kiểm tra tổng khối lượng các kiện hàng của một hoạt động không vượt quá \texttt{drone.max\_payload\_kg} cần phép tổng hợp trên nhiều bản ghi kiện hàng.
        \item \textbf{Tính đồng thời của chuyến bay (BR-16)}: Ngăn lịch bay chồng lấn của cùng một máy bay khi đang hoạt động cần ràng buộc loại trừ theo thời gian (\texttt{EXCLUDE USING gist}) hoặc trigger kiểm tra trước khi kích hoạt hoạt động.
        \item \textbf{Vòng đời chuyển trạng thái (BR-12, BR-13, BR-14)}: Ngăn chuyển lùi trạng thái không hợp lệ, chẳng hạn từ \texttt{DELIVERED} về \texttt{ASSIGNED}, cần trigger kiểm tra chuyển trạng thái hoặc máy trạng thái tại tầng ứng dụng.
        \item \textbf{Quyền của nhân viên đang được phân công (BR-28)}: Xác nhận người dùng ghi sự kiện kiện hàng đang được phân công vào trạm đó cần trigger đối chiếu với \texttt{station\_operator\_assignment}.
    \end{itemize}
\end{enumerate}
Việc ghi rõ những ranh giới này bảo đảm tính trung thực học thuật và tránh các tuyên bố thiếu cơ sở khi bảo vệ dự án.

\section{Tóm tắt chương}\label{s:ch4summary}

\qquad Chương này đã kiểm chứng có hệ thống thiết kế cơ sở dữ liệu SmartDroneDelivery được đề xuất. Mọi khía cạnh cấu trúc, bao gồm khóa chính, khóa ngoại, khóa ứng viên, ràng buộc kiểm tra miền và hành động tham chiếu, đều được đối chiếu với các yêu cầu hệ thống.

Ma trận kiểm chứng xác nhận cả ba mươi chín quy tắc nghiệp vụ đều được biểu diễn rõ ràng trong mô hình quan hệ. Các kỹ thuật quan hệ nâng cao như khóa ngoại ghép để bảo đảm quyền sở hữu địa chỉ khách hàng (BR-04) và chỉ mục duy nhất từng phần cho địa chỉ mặc định (BR-05) đã được minh họa. Sự tách biệt về cấu trúc giữa lịch sử trạng thái nghiệp vụ và nhật ký kiểm toán bảo mật đã được kiểm chứng, đồng thời việc tuân thủ 3NF cũng được xác nhận.

Cuối cùng, chương đã phân định rõ ràng ranh giới giữa các ràng buộc cơ sở dữ liệu khai báo và logic nghiệp vụ ở cấp thủ tục/ứng dụng, qua đó cung cấp một nền tảng trung thực, chặt chẽ và có thể bảo vệ cho thiết kế cơ sở dữ liệu.
